\section{本讲小结}
极值分析要从定义域、候选点和导数符号展开，势能判稳则必须明确允许的扰动方向。近似展开先找无量纲小参数，再保留足以决定结果的最低非零阶；发生相减时应特别检查低阶项是否抵消。曲率半径描述路径的局部几何，几何点的速度则还取决于驱动条件，二者需要不同的信息。
\section{练习}
\begin{question}
\begin{drill}[带误差的近似]\label{c3d1}
用二阶展开估算 $\sqrt{4.01}$，并给出误差上界。再估计把 $\sin\theta$ 近似为 $\theta$ 时，在 $|\theta|\le0.1$ 内相对 $\theta$ 的误差上界。
\end{drill}
\fig{c3d1q}{展开点与实际求值点相邻}{\begin{tikzpicture}
\draw[axis](0,0)--(5.5,0)node[right]{$x$};\node[below]at(1,0){$4$};\node[below]at(4,0){$4.01$};\draw[dim](1,.5)--(4,.5)node[midway,above]{$h=0.01$};
\end{tikzpicture}}\answerlink{c3d1}
\end{question}
\begin{question}
\begin{drill}[分层大气中的圆形光路]\label{c3d2}
球对称大气折射率为 $n(r)=1+\delta\e^{-(r-r_0)/b}$，$r_0=6400\,\mathrm{km}$，$b=8772\,\mathrm m$。已知近地层实际折射率为 $1.0003$。若光线要沿 $r=r_0$ 的圆周传播，所需密度约为实际密度的多少倍？假定 $n-1$ 与密度成正比，并使用分层介质光线曲率关系 $\kappa=-n'(r)/n(r)$。
\end{drill}
\fig{c3d2q}{折射率梯度朝向地心}{\begin{tikzpicture}[scale=.9]
\draw[world](0,0)circle(1.4);\draw[guide](0,0)circle(1.8);\draw[ray,blue](1.8,0)arc(0:95:1.8);
\draw[flow](40:2.2)--(40:1.5)node[right]{$\nabla n$};\node[below]at(0,0){$O$};\draw[guide](0,0)--(0,1.8)node[midway,right]{$r_0$};
\end{tikzpicture}}\answerlink{c3d2}
\end{question}
\begin{question}
\begin{drill}[闭区间的最大与最小]\label{c3d3}
求 $f(x)=x^3-3x$ 在 $[-2,3]$ 上的最大值和最小值，说明为何只解 $f'=0$ 不足以完成此题。
\end{drill}
\fig{c3d3q}{端点也属于可取位置}{\begin{tikzpicture}
\draw[axis](-2.7,0)--(3.7,0)node[right]{$x$};\draw[very thick,blue](-2,0)--(3,0);\foreach \x in {-2,-1,1,3}{\fill(\x,0)circle(2pt)node[below]{$\x$};}
\end{tikzpicture}}\answerlink{c3d3}
\end{question}
\begin{question}
\begin{drill}[曲率与参数选择]\label{c3d4}
质点轨迹为 $x=A\cos\omega t,\ y=B\sin\omega t$，$A>B>0,\omega>0$。求 $t=0$ 时速率、加速度和曲率半径；再求 $\omega t=\pi/4$ 时的切向加速度，判断它是否一般为零。
\end{drill}
\fig{c3d4q}{恒定参数角速度不等于恒定速率}{\begin{tikzpicture}
\draw[world](0,0)ellipse(2 and 1);\fill(2,0)circle(2pt)node[right]{$t=0$};\fill(1.4142,.7071)circle(2pt)node[above right]{$\omega t=\pi/4$};
\draw[guide](0,0)--(1.4142,.7071);\node[below]at(0,0){$O$};
\end{tikzpicture}}\answerlink{c3d4}
\end{question}
